\BourbakiExerciseGroup

\begin{enumerate}
\def\labelenumi{\arabic{enumi})}
\item
设 \(f\) 是开区间 \(\mathbf{I} \subset \mathbb{R}\) 到 \(\mathbb{R}\) 中的一个连续映射。试证：若 \(f(\mathbf{I})\) 是开集，且对每个 \(y \in \mathbb{R}\) ，集合 \(\overline{f}^1(y)\) 至多含两个不同的点，则 \(f\) 是单调的。
\item
设 B 是把 R 看作域 Q 上的向量空间时的一个基（``Hamel 基''）。B 是不可数的（§ 2，习题 2）。设 φ 是 B 的一个子集 C ≠ B 到集合 B 上的一个双射。如下定义 R 到 R 中的一个映射 f：若 \(x = \sum_{\xi \in B} \lambda(\xi) \xi\) ，其中 \(\lambda(\xi) \in \mathbf{Q}\) 且 \(\lambda(\xi) = 0\) 对除有限个元素 \(\xi \in B\) 外的一切元素成立，则 \(f(x) = \sum_{\xi \in C} \lambda(\xi) \varphi(\xi)\) 。试证：\(f(x + y) = f(x) + f(y)\) ，但对每个 \(z \in R\) ， \(f(z)\) 是 R 的一个稠密子集，由此推出 f 在 R 的任何区间上既不上有界也不下有界。
\item
  \begin{enumerate}
  \def\labelenumii{\alph{enumii})}
  \tightlist
  \item
对每个区间 \(\mathbf{I} \subset \mathbb{R}\) ，用 \(\mathbf{G}(\mathbf{I})\) 表示 \(\mathbf{I}\) 到其自身的所有同胚所成的群。试证：若 \(\mathbf{I}\) 与 \(\mathbf{J}\) 是 \(\mathbb{R}\) 的任意两个区间，每个都含多于一个点，则 \(\mathbf{G}(\mathbf{I})\) 与 \(\mathbf{G}(\mathbf{J})\) 同构。在 \(\mathbf{G}(\mathbf{I})\) 中，递增同胚所成的集 \(\mathbf{F}(\mathbf{I})\) 是指数为 2 的正规子群。
  \end{enumerate}
\end{enumerate}

\begin{enumerate}
\def\labelenumi{\alph{enumi})}
\setcounter{enumi}{1}
\item
设 \(\mathbf{G} = \mathbf{G}(\mathbf{R})\) ， \(\mathbf{F} = \mathbf{F}(\mathbf{R})\) 。对每个 \(f \in \mathbf{G}\) 与每个 \(x \in \mathbf{R}\) ，用 \(\sigma(x; f)\) 表示 \(\operatorname{sgn}(f(x) - x)\) 。设 \(\mathbf{T}^{+}\) （相应地 \(\mathbf{T}^{-}\) ）是所有这样的 \(f \in \mathbf{F}\) 所成之集：\(\sigma(x; f)\) 等于 \(\mathfrak{I}\) （相应地 \(-\mathfrak{I}\) ）且在 \(\mathbf{R}\) 上为常数。试证：\(\mathbf{T}^{+}\) （相应地 \(\mathbf{T}^{-}\) ）的每个元素在 \(\mathbf{F}\) 中共轭于平移 \(x \to x + \mathfrak{I}\) （相应地 \(x \to x - \mathfrak{I}\) ）。{[}若 \(f \in \mathbf{T}^{+}\) ，考察点列 \(f^{n}(0), n \in \mathbf{Z}\) {]}。
\item
设 \(T = T^{+} \cup T^{-}\) 。试证：每个不属于 T 的 \(f \in F\) 都是 \(T_{*}\) 中两个元素在 F 中的乘积{[}在 F 中写 \(f = f_{1}f_{2}\) ，其中 \(f_{1}(x) = \sup (x, f(x))\) ， \(f_{2}(x) = \inf (x, f(x))\) ；若 g 是平移 \(x \to x + 1\) ，则 \(f_{1}g\) 与 \(g^{-1}f_{2}\) 属于 T{]}。
\item
两个元素 \(f, g\) 在 \(F\) 中共轭，当且仅当存在 \(s \in F\) ，使得 \(\sigma(x; f) = \sigma(s(x); g)\) 对一切 \(x \in \mathbb{R}\) 成立。{[}考察连通分支 \(I_n\) ——即所有 \(x\) （满足 \(\sigma(x; f) \neq 0\) ）所成的开集的连通分支；注意 \(F(I_n)\) 与 \(F\) 同构（由 \(a\) ）并用 \(b\) {]}。特别地，若 \([a, b]\) 是 \(R\) 的一个区间，使得 \(\sigma(a; f) = \sigma(b; f) = 0\) 对某个 \(f \in F\) 成立，则如下元素 \(f'\) 属于 \(F\) ：它与 \(f\) 相重合（对 \(x \leqslant a\) 或 \(x \geqslant b\) ），且满足 \(f'(x) = x + \varepsilon(f(x) - x)\) （对 \(a < x < b\) ，其中 \(0 < \varepsilon < 1\) ）——这个元素与 \(f\) 在 \(F\) 中共轭。
\item
设 \(f \in F\) ，并设 \(a, b, c, d\) 是四个实数，满足 \(a < c < b < d\) ，且 \(f(x) - x = 0\) 对 \(x = a, x = b, x = c, x = d, \sigma(x; f) = 1\) 成立（于 \(a < x < c\) 与 \(b < x < d\) ）。设 \(a', b'\) 是两个数，满足 \(a \leqslant a' < c\) 与 \(b < b' \leqslant d\) ；令 \(J = [a, b], J' = [a', b']\) ，并用 \(f_1\) 表示 \(f\) 在 \(J\) 上的限制；于是 \(f_1 \in F(J)\) 。试证：存在一个递增同胚 \(s\) （ \(J\) 到 \(J'\) 上的），使得元素 \(f_2 = sf_1s^{-1}\) （ \(F(J')\) 的）满足：\(f_2(x) > f^{-1}(x)\) 每当 \(a' < x < b'\) {[}用 \(d\) {]}。
\end{enumerate}

\(f)\) 设 \(f \in \mathbf{F}\) ，并设 \([a, b]\) 是 \(\mathbf{R}\) 的一个区间，使得 \(f(x) - x = 0\) 对 \(x = a\) 与 \(x = b\) 成立，\(\sigma(x; f) = -1\) 对 \(x < a\) 成立，且 \(\sigma(x; f) = +1\) 对 \(x > b\) 成立。试证：存在 \(g \in \mathbf{F}\) ，它共轭于 \(f\) （在 \(\mathbf{F}\) 中），使得 \(g(x) > f(x)\) 对一切 \(x \in \mathbf{R}\) 成立（同样的方法）。

\(g\) ) 用 \(\mathbf{H}^{+}\) （相应地 \(\mathbf{H}^{-}\) ）表示 \(\mathbf{F}\) 的由所有这样的 \(f\in \mathbf{F}\) 组成的正规子群：\(f(x) = x\) 在 \(+\infty\) （相应地 \(-\infty\) ）的一个邻域内成立。试证：\(\mathbf{H}^{+},\mathbf{H}^{-}\) 与 \(\mathbf{H} = \mathbf{H}^{+}\cap \mathbf{H}^{-}\) 是 \(\mathbf{F}\) 的除 \(\mathbf{F}\) 与 \(\{e\}\) 外仅有的正规子群。{[}若 \(\mathbf{N}\) 是 \(\mathbf{F}\) 的一个正规子群，异于 \(\mathbf{F}\) ，且 \(f\in \mathbf{N}\) 不属于 \(\mathbf{H}^{+}\) ，试证存在一个元素 \(g\in \mathbf{N}\) ，使得所有 \(x\in \mathbf{R}\) （对其有 \(\sigma (x;g) = +1\) ）所成之集是一个区间 \([a, + \infty ]\) ，并且 \(g(x) = x\) 对 \(x\leqslant a\) 成立。为此，利用 \(e)\) 与 \(f)\) 的构造，以及 \(b)\) 与 \(c)\) 。要证明 \(\mathbf{N}\) 不包含除 \(\mathbf{H}\) 与 \(\{e\}\) 外的在 \(\mathbf{F}\) 中正规的子群，考察一个元素 \(f\in \mathbf{H}\) ，它异于 \(e\) ；设 \(a\) 与 \(b\) 是所有这样的 \(x\in \mathbf{R}\) （\(\sigma (x;f)\neq 0\) ）所成之集的下确界与上确界（由假设二者均有限）。然后注意 \(\mathrm{F}([a,b])\) 同构于 \(\mathbf{F}\) ，并用前面的结果{]}。

\begin{enumerate}
\def\labelenumi{\alph{enumi})}
\setcounter{enumi}{7}
\tightlist
\item
试证：群 H 是单群。（与证明 H 不含在 F 中正规的非平凡子群的方法相同）。
\end{enumerate}

\begin{enumerate}
\def\labelenumi{\arabic{enumi})}
\setcounter{enumi}{3}
\item
设 \(f\) 是拓扑空间 \(\mathbf{X}\) 上的一个下半连续函数，\(\varphi\) 是 \(f(\mathbf{E}) \subset \overline{\mathbf{R}}\) 上的一个递增下半连续函数。试证：\(\varphi \circ f\) 在 \(\mathbf{X}\) 上是下半连续的。
\item
设 \(f\) 是拓扑空间 \(\mathbf{X}\) 上的一个下半连续函数。试证：若 \(\mathbf{A}\) 是 \(\mathbf{X}\) 的任一非空子集，则 \(\sup f(\overline{\mathbf{A}}) = \sup f(\mathbf{A})\) 。
\item
设 \(f\) 是定义在拓扑空间 \(X\) 上的一个实值函数。数（有限或无穷）
\end{enumerate}

\[
\omega (a; f) = \lim _ {x \to a} \sup f (x) - \lim _ {x \to a} \inf f (x)
\]

称为 \(f\) 在点 \(a \in \mathbf{X}\) 处的振荡，只要右边有定义。

\begin{enumerate}
\def\labelenumi{\alph{enumi})}
\item
试证：\(x \to \omega(x; f)\) 在子空间 \(A\) （ \(X\) 的）上是上半连续的，该映射正是在其上有定义。
\item
若 \(f\) 在 \(\mathbf{X}\) 上有限，则 \(\mathbf{A} = \mathbf{X}\) ，且对每个 \(a \in \mathbf{X}\) 有
\end{enumerate}

\[
\omega (a; f) = \lim _ {(x, y) \rightarrow (a, a)} \sup (f (x) - f (y)).
\]

\begin{enumerate}
\def\labelenumi{\alph{enumi})}
\setcounter{enumi}{2}
\tightlist
\item
设 \(f\) 是 \(\mathbf{X}\) 上的一个下半连续有限实值函数。试证：若在一处 \(a \in \mathbf{X}, \omega(a; f)\) 取有限值，则
\end{enumerate}

\[
\liminf _ {x \to a} \omega (x; f) = 0.
\]

{[}设结论不真，试证将存在任意接近 a 的点 x，使 \(f(x)\) 可以任意大{]}。

\begin{enumerate}
\def\labelenumi{\alph{enumi})}
\setcounter{enumi}{3}
\tightlist
\item
对每个既约形式的有理数 \(r = p / q\) （其中 \(q > 0\) ），令 \(f(r) = q\) 。试证：\(f\) 在 \(\mathbf{Q}\) 上是下半连续的，但在每个点 \(r \in \mathbf{Q}\) 处有 \(\omega(r; f) = +\infty\) 。
\end{enumerate}

¶ 7) 设 X 是局部紧空间，f 是 X 上的一个下半连续函数。

\begin{enumerate}
\def\labelenumi{\alph{enumi})}
\tightlist
\item
若 \(\limsup_{x\to a}f(x)=+\infty\) 对每个 \(a\in\mathbf{X}\) 成立，试证：集 \(\overline{f}(+\infty)\) 在 \(\mathbf{X}\) 中稠密{[}对每个 \(a\in\mathbf{X}\) 与每个邻域 \(\mathbf{V}\) （ \(a\) 的），定义一个相对紧开集序列 \((\mathbf{U}_{n})\) ，使得 \(\overline{\mathbf{U}}_{n+1}\subset\mathbf{U}_{n}\) ，且 \(f(x)>n\) 对一切 \(x\in\mathbf{U}_{n}\) 成立{]}。
\end{enumerate}

* b) 设 \(n \to r_n\) 是 \(\mathbf{N}\) 到属于 {[}0, 1{]} 的有理数所成之集上的一个双射，并令 \(\varphi(x) = \mathrm{i}/\sqrt{|x|}\) {[}取 \(\varphi(0) = +\infty\) {]}；则函数 \(f(x) = \sum_{n=0}^{\infty} 2^{-n} \varphi(x - r_n)\) 在 {[}0, 1{]} 上是下半连续的，\(\overline{f}(+\infty)\) 在这个区间中稠密，其余集也如此{[}为证最后一点，注意以

\[
2 ^ {- n} \int_ {0} ^ {1} \varphi (x - r _ {n}) d x. *
\]

（积分学，第四章，§ 3，第 6 节，定理 5）。{]} *

\begin{enumerate}
\def\labelenumi{\alph{enumi})}
\setcounter{enumi}{2}
\item
若 \(f\) 有限，试证：点 \(x\) （使 \(\omega(x; f)\) 有限）所成之集是稠密的{[}用 a){]}。
\item
试证：X 中使 f 连续的点（无论 f 有限与否）所成之集是稠密的。{[}用 ( f/(1 + \textbar f\textbar) ) 代替 f（习题 4），化归到 f 有界的情形；注意 \(1/\omega(x; f)\) 是下半连续的，并用 a) 与习题 6c){]}。
\end{enumerate}

\begin{enumerate}
\def\labelenumi{\arabic{enumi})}
\setcounter{enumi}{7}
\item
设 X 是拓扑空间，A 是 \(X \times \overline{R}\) 的一个闭子集。试证：映射 \(x \to \inf(A(x))\) （ \(pr_{1}A\) 到 \(\overline{R}\) 中的）是下半连续的。反之，若 \(f: X \to \overline{R}\) 是下半连续函数，则 \(X \times \overline{R}\) 的子集 B（由对 \((x, y)\) （满足 \(y \geqslant f(x)\) ）组成）在 \(X \times \overline{R}\) 中是闭的。
\item
  \begin{enumerate}
  \def\labelenumii{\alph{enumii})}
  \tightlist
  \item
设 X 与 Y 是两个分离空间，\(\pi: X \to Y\) 是一个逆紧映射（第一章，§ 10，第 1 节）。设 \(g\) 是 X 上的一个下半连续实值函数。对每个 \(y \in Y\) ，令 \(f(y)\) 为 \(g\) 在集合 \(\overline{\pi}^{1}(y)\) 中的下确界{[}于是 \(f(y) = +\infty\) 若 \(\overline{\pi}^{1}(y) = \emptyset\) {]}。试证：\(f\) 在 Y 上是下半连续的。{[}利用如下事实：\(\pi(X)\) 是闭集，\(\overline{\pi}^{1}(y)\) 对一切 \(y \in \pi(X)\) 是紧的，以及 \(\overline{\pi}^{1}(y)\) 的每个邻域都包含一个关于等价关系 \(\pi(x) = \pi(x')\) 饱和的邻域。{]}
  \end{enumerate}
\end{enumerate}

\begin{enumerate}
\def\labelenumi{\alph{enumi})}
\setcounter{enumi}{1}
\tightlist
\item
设 \(g\) 是连续函数 \(|x_1 x_2 - 1|\) （定义在 \(\mathbf{R} \times ]0, +\infty[\) 上），对每个 \(x_1 \in \mathbb{R}\) 令 \(f(x_1) = \inf_{x_1 > 0} g(x_1, x_2)\) 。试证：\(f\) 不是下半连续的。
\end{enumerate}

\begin{enumerate}
\def\labelenumi{\arabic{enumi})}
\setcounter{enumi}{9}
\tightlist
\item
把定义 1、命题 1 与定理 3 推广到定义在拓扑空间 X 上、取值于任意全序集 Y 中的函数。若 f 与 g 是 X 到 Y 中的两个映射，都在一点 a 处下半连续，则 inf (f, g) 与 sup (f, g) 在 a 处下半连续。若 Y 是全序交换群，则 \(f + g\) 亦然。
\end{enumerate}

若 \(\mathbf{Y}\) 在拓扑 \(\mathfrak{G}_{0}(\mathbf{Y})\) （ \(\S 2\) ，习题 6）下是紧的，把定理 4 与命题 3、4 推广到取值于 \(\mathbf{Y}\) 中的函数。

¶ 11) 设 \(\overline{\mathbf{R}}\) 带有右拓扑（第一章，§ 1，习题 2）。一个映射 \(f\) （拓扑空间 \(\mathbf{X}\) 到 \(\overline{\mathbf{R}}\) 中的）在这个拓扑下连续，当且仅当 \(f\) 在 \(\mathbf{X}\) 的每一点处有相对极小。若 \(\mathbf{X} = \mathbf{R}\) （带通常拓扑），试证：集合 \(f(\mathbf{X})\) 于是是可数的。{[}设结论不真，且区间 \([\alpha, \beta] \cap f(\mathbf{X})\) 不可数；对每个 \(y \in [\alpha, \beta]\) ，令 \(\mathbf{U}(y) = \overline{f}([y, +\infty])\) ；\(\mathbf{U}(y)\) 是 \(\mathbf{R}\) 中的开子集。设 \((\mathbf{I}_n)\) 是 \(\mathbf{U}(\beta)\) 的连通分支的（有限或无穷）序列，对每个 \(n\) 令 \(x_n\) 为 \(\mathbf{I}_n\) 的一个点；对每个 \(y \in [\alpha, \beta]\) ，令 \(g_n(y)\) （相应地 \(h_n(y) \) ）为 \(\mathbf{U}(y)\) 的包含 \(x_n\) 的那个连通分支的左（相应地右）端点。试证：对至少一个 \(n\) 的值，函数 \(g_n, h_n\) 之一必在 \([\alpha, \beta]\) 中取不可数无穷多个不同的值；然后用 § 5 习题 8{]}。

\begin{enumerate}
\def\labelenumi{\arabic{enumi})}
\setcounter{enumi}{11}
\item
设 \(f\) 是紧区间 \([a, b]\) （ \(\mathbf{R}\) 的）上的一个连续、非常数的有限实值函数，满足 \(f(a) = f(b) = 0\) ；设 \(Z\) 是 \([a, b]\) 中由点 \(x\) （对其有 \(f(x) = 0\) ）组成的闭子集，\(c\) 是邻接于 \(Z\) 的诸区间在 \([a, b]\) 中长度的最大者。试证：对每个 \(t\) （满足 \(0 < t < c\) ），存在一点 \(x \in [a, b]\) ，使得 \(x + t \in [a, b]\) 且 \(f(x + t) = f(x)\) 。
\item
设 \(f\) 是紧区间 \([a, b]\) （ \(\mathbf{R}\) 的）上的一个连续有限实值函数，并设 \(f\) 不是严格单调的。试证：存在一点 \(x_0 \in ]a, b[\) ，使得对每个 \(\varepsilon > 0\) ，有两点 \(y, z\) 在 \([a, b]\) 中满足 \(x_0 - \varepsilon < y < x_0 < z < x_0 + \varepsilon\) 且 \(f(y) = f(z)\) 。
\item
设 \(f\) 是 \(\mathbf{R}\) 到其自身的一个连续映射。
\end{enumerate}

\begin{enumerate}
\def\labelenumi{\alph{enumi})}
\tightlist
\item
试证：若 \(f\) 在 \(\mathbf{R}\) 上一致连续，则存在两个实数 \(\alpha \geqslant 0\) 与 \(\beta \geqslant 0\) ，使得 \(|f(x)| \leqslant \alpha |x| + \beta\) 对一切 \(x \in \mathbf{R}\) 成立。b) 试证：若 \(f\) 在 \(\mathbf{R}\) 中单调且有界，则 \(f\) 在 \(\mathbf{R}\) 上一致连续。
\end{enumerate}

¶ 15) 设 X 是 Alexandroff 半直线{[}§ 2，习题 12 b){]}，并设 X′ 是它通过添加一个无穷远点 ω 所得的紧化。试证：对每个连续映射 f: X → X，或者我们有

\(\lim_{x \to \omega, x \in \mathbf{X}} f(x) = \omega,\) 或者存在 \(z \in \mathbf{X}\) 使得 \(f\) 为常数

在区间 \([z, \rightarrow[\) 中。{[}试证：若当 \(x\) 趋于 \(\omega, f\) 时有一个聚点 \(c \in \mathbf{X}\) ，则必有 \(\lim_{x \to \omega, x \in \mathbf{X}} = c\) ；证明 \(f\) 在 X′ 中不能有其他聚点，利用每个递增序列在 X 中收敛这一事实。然后用 § 2 习题 12 b) 以及在 X′ 中 ω 的任何可数个邻域之交仍是 ω 的邻域这一事实{]}。

